\documentclass{article}

\usepackage{amsmath,amssymb}
\usepackage{xurl}
\usepackage[hidelinks]{hyperref}

% Shared commands for notes in docs/paper-gaps/.

\DeclareUnicodeCharacter{2080}{\ifmmode{}_{0}\else\textsubscript{0}\fi}
\DeclareUnicodeCharacter{2081}{\ifmmode{}_{1}\else\textsubscript{1}\fi}
\DeclareUnicodeCharacter{2082}{\ifmmode{}_{2}\else\textsubscript{2}\fi}
\DeclareUnicodeCharacter{2099}{\ifmmode{}_{n}\else\textsubscript{n}\fi}

\newcommand{\Lean}{\textsc{Lean}}
\ExplSyntaxOn
\NewDocumentCommand{\leanid}{m}{
  \ifmmode
    \textsf{\detokenize{#1}}
  \else
    \group_begin:
    \urlstyle{sf}
    \tl_set:Nn \l_tmpa_tl {#1}
    \tl_replace_all:Nnn \l_tmpa_tl {\_} {_}
    \exp_args:NV \path \l_tmpa_tl
    \group_end:
  \fi
}
\ExplSyntaxOff
\providecommand{\tr}{\operatorname{tr}}
\newcommand{\tnleanIssueBase}{https://github.com/LionSR/TNLean/issues/}
\newcommand{\tnleanPrBase}{https://github.com/LionSR/TNLean/pull/}
\newcommand{\tnleanDocBase}{https://sirui-lu.com/TNLean/docs/}
\newcommand{\ghissue}[1]{\href{\tnleanIssueBase#1}{GitHub issue \##1}}
\newcommand{\ghpr}[1]{\href{\tnleanPrBase#1}{PR \##1}}
\newcommand{\doclink}[1]{\href{\tnleanDocBase#1.html}{\path{#1}}}

% Dirac notation and the identity operator, as in blueprint/src/macros/common.tex.
% \providecommand keeps note-local definitions authoritative.
\usepackage{dsfont}
\providecommand{\ket}[1]{|#1\rangle}
\providecommand{\bra}[1]{\langle #1|}
\providecommand{\braket}[2]{\langle #1 | #2 \rangle}
\providecommand{\ketbra}[2]{|#1\rangle\!\langle #2|}
\providecommand{\Id}{\mathds{1}}
\providecommand{\one}{\mathds{1}}
\providecommand{\C}{\mathbb{C}}
\providecommand{\MN}[1]{M_{#1}(\C)}
\providecommand{\MD}{M_D(\C)}

% External referencing for paper sources.  The build helper writes sanitized
% auxiliary files under build/paper-aux/ and excludes duplicated source labels.
\usepackage{xr}

% Import a generated paper auxiliary file with a caller-supplied label prefix.
% The candidate paths support builds from docs/paper-gaps/, the repository root,
% and build/paper-aux/ itself.
\newcommand{\importpaperaux}[2]{%
  \IfFileExists{../../build/paper-aux/#2.aux}{%
    \externaldocument[#1]{../../build/paper-aux/#2}%
  }{%
    \IfFileExists{build/paper-aux/#2.aux}{%
      \externaldocument[#1]{build/paper-aux/#2}%
    }{%
      \IfFileExists{#2.aux}{%
        \externaldocument[#1]{#2}%
      }{}%
    }%
  }%
}

% General external reference macros
\newcommand{\paperref}[2]{\ref{#1-#2}}
\newcommand{\papereqref}[2]{(\ref{#1-#2})}

% CPSV16 source (1606.00608 / MPDO-22-12-17-2.tex) external document and shortcuts
\importpaperaux{cpsv16-}{1606.00608/MPDO-22-12-17-2-tnlean}
\newcommand{\cpsvref}[1]{\paperref{cpsv16}{#1}}
\newcommand{\cpsveqref}[1]{\papereqref{cpsv16}{#1}}

% Verdict marker: \gapnote{<kind>}{<status>}. Kind is the policy
% classification (clarification, local-correction, scope-restriction,
% unfaithful, false-source, open-gap); status is open, wip (elimination
% underway), resolved, or historical. Machine-read by the site index;
% prints a verdict line.
\newcommand{\gapnote}[2]{\par\noindent\textbf{Verdict:} #1 (#2).\par}

\importpaperaux{fbc25-}{2502.20257/main-tnlean}

\title{CZX Defect Domains from the Source Diagrams}
\author{}
\date{2026-09-04}

\begin{document}
\maketitle
\gapnote{clarification}{open}

\section{The assertion}

The state-level gauging of the anomalous $\mathbb Z_2$ CZX symmetry in
Ref.~\cite[lines~4503--5183]{FrancoRubio2025MPUGauging} rests on the
modified-fusion lemma of
Ref.~\cite[lines~4215--4254]{FrancoRubio2025MPUGauging}. That lemma attaches
to every ordered pair of group labels $(g,h)$ and every injective block $x$ a
two-site defect vector $v(g,h,x)$, drawn as the contraction of the defect
tensor $g[hx]$ on the left site with the defect tensor $h[x]$ on the right
site, and prescribes unitary fusion operators $\tilde\lambda_{g,h}$ on the
two-site matter space with
\begin{align}
    \tilde\lambda_{g,h}\,v(g,h,x)
        &=v(e,gh,x).
    \label{eq:fbc25_czx_prescription}
\end{align}
For the CZX example the source displays the blocked symmetry tensor, its two
decompositions, the two injective blocks of the invariant GHZ state, the
action tensors, the two domain-wall defect tensors, the movement operator,
the fusion operators, and the modification of the $(g,g)$ fusion operator.
It does not display the eight two-site vectors $v(g,h,x)$ themselves, nor
the four defect domains
\begin{align}
    \mathcal K_{g,h}
        &=\operatorname{span}\{v(g,h,x):x\in\{0,1\}\}
    \label{eq:fbc25_czx_domain_definition}
\end{align}
on which Eq.~\ref{eq:fbc25_czx_prescription} constrains the fusion
operators. Unpublished planning material for the gauge-invariant-subspace
problem therefore treated two statements as explicit axioms: that the
retained four-qubit circuits coincide with the source's fusion operators in
one common blocking, basis, phase, and fusion gauge, and that the odd target
support $\operatorname{span}\{v(e,g,x):x\}$ coincides with the even one
$\operatorname{span}\{v(e,e,x):x\}$. This note derives all eight vectors by
finite contraction of the displayed tensors and settles both statements.
The first is a theorem. The second is false: the two target supports are
orthogonal.

\section{Conventions read off the diagrams}

One blocked site carries two qubits. Two neighboring blocked sites carry the
qubits $0,1$ (left site) and $2,3$ (right site), and the computational basis
of the sixteen-dimensional two-site matter space is written
$\ket{x_0x_1x_2x_3}$. The two injective blocks of the once-blocked GHZ state
are the product vectors
\begin{align}
    \ket{\circ}&=\ket{00},
    &
    \ket{\bullet}&=\ket{11},
    \label{eq:fbc25_czx_blocks}
\end{align}
drawn as a white and a black circle
\cite[lines~4660--4670]{FrancoRubio2025MPUGauging}. Both blocks have bond
dimension one, so every virtual contraction between neighboring defect
tensors is a scalar and every single-site defect tensor is a two-qubit
vector.

Circuit diagrams are read from the bottom (input) to the top (output), and
operator products act from right to left. The controlled-$Z$ gate is drawn
as a bond of dimension two between the two qubit lines with a blue dot on the
bond \cite[lines~3800--3845]{FrancoRubio2025MPUGauging}: the attachment on a
qubit line is the copy tensor $\delta_{i,j,k}$, which exports the qubit value
onto the bond, and the blue dot is the matrix
\begin{align}
    \sqrt2H
        &=\begin{pmatrix}1&1\\1&-1\end{pmatrix},
    &
    \mathrm{CZ}_{rs}\ket{x}
        &=(-1)^{x_rx_s}\ket{x}.
    \label{eq:fbc25_czx_cz_decomposition}
\end{align}
The blocked symmetry tensor \papereqref{fbc25}{eq:MPU_CZX} of
Ref.~\cite[lines~4505--4545]{FrancoRubio2025MPUGauging} applies $X\otimes X$
at the bottom, then a controlled-$Z$ between its two qubits, exports the
first qubit onto the left bond and the second qubit onto the right bond
(with the blue dot on the right bond), and applies $Z\otimes Z$ at the top.
The source picks two decompositions of this tensor
\cite[lines~4559--4659]{FrancoRubio2025MPUGauging}. The one used in the
definition of the defect tensors is the white tensor with a wavy output leg
and a bond leg to the left: it applies $Z\otimes Z$, exports the first qubit
onto the bond through a copy tensor, and carries the Pauli matrix $Y$ on the
bond. On a block vector $\ket{xx}$ its output is therefore
\begin{align}
    \ket{xx}\otimes Y\ket{x}
        &=i(-1)^{x}\,\ket{xx}\otimes\ket{x+1},
    \label{eq:fbc25_czx_white_tensor_on_block}
\end{align}
where the second factor lives on the bond. The truncated symmetry displayed
in Ref.~\cite[lines~4700--4798]{FrancoRubio2025MPUGauging}, whose rightmost
site carries this white tensor, has the overall prefactor $i$ that
Eq.~\ref{eq:fbc25_czx_white_tensor_on_block} produces; this confirms the
direction in which $Y$ acts on the bond. The left action tensors are the
bond covectors
\begin{align}
    A^{\Gamma}_{g,0}&=-i\bra{1},
    &
    A^{\Gamma}_{g,1}&=\bra{0}
    \label{eq:fbc25_czx_action_tensors}
\end{align}
of Ref.~\cite[lines~4690--4694]{FrancoRubio2025MPUGauging}, labelled by the
block $x$ to the right of the defect.

\section{The single-site defect vectors}

The domain-wall defect $g[x]$ is defined in
Ref.~\cite[lines~2831--2845]{FrancoRubio2025MPUGauging} as the block tensor
$x$ at the bottom, the white tensor of
Eq.~\ref{eq:fbc25_czx_white_tensor_on_block} on its physical leg, and the
left action tensor contracted against the bond. Since $g$ exchanges the two
blocks, the defect $g[1]$ carries block $0$ to its left and block $1$ to its
right, which the source writes $g[1]\equiv0|1$, and similarly
$g[0]\equiv1|0$. Contracting Eq.~\ref{eq:fbc25_czx_white_tensor_on_block}
with Eq.~\ref{eq:fbc25_czx_action_tensors} gives
\begin{align}
    g[1]
        &=A^{\Gamma}_{g,1}\bigl(i(-1)^{1}\ket{0}\bigr)\,\ket{11}
        =-i\ket{11},
    \label{eq:fbc25_czx_defect_g1}\\
    g[0]
        &=A^{\Gamma}_{g,0}\bigl(i\ket{1}\bigr)\,\ket{00}
        =\ket{00}.
    \label{eq:fbc25_czx_defect_g0}
\end{align}
These are exactly the evaluations displayed in
Ref.~\cite[lines~4800--4859]{FrancoRubio2025MPUGauging}: $g[1]$ equals $-i$
times the black circle and $g[0]$ equals the white circle. The displayed
values fix the reading of the diagram. The transposed action of $Y$ on the
bond would give $g[1]=i\ket{11}$ and $g[0]=-\ket{00}$, contradicting both
displayed evaluations. The alternative expressions displayed in the same
lines, through the black tensor of the first decomposition and the right
action tensors $-\tfrac{i}{\sqrt2}\ket{-}$ for block $0$ and
$\tfrac1{\sqrt2}\ket{+}$ for block $1$, labelled by the block to the left of
the defect, reproduce Eqs.~\ref{eq:fbc25_czx_defect_g1}
and~\ref{eq:fbc25_czx_defect_g0} with the displayed signs, that is, with
the sign function $\xi_g(1)=-1$ and $\xi_g(0)=1$ of
Ref.~\cite[lines~2925--2962]{FrancoRubio2025MPUGauging}.

The identity defect $e[x]$ is the block tensor itself,
\begin{align}
    e[0]&=\ket{00},
    &
    e[1]&=\ket{11}.
    \label{eq:fbc25_czx_defect_e}
\end{align}
This is the source's convention: projecting the gauge degrees of freedom of
the gauged tensor onto the neutral element recovers the original tensor
\cite[lines~3494--3529]{FrancoRubio2025MPUGauging}, and every displayed
fusion identity in
Ref.~\cite[lines~4889--5066]{FrancoRubio2025MPUGauging} draws a fused
identity defect as a plain block circle.

\section{The two-site defect vectors, domains, and prescribed maps}

The two-site vector $v(a,b,x)$ is the tensor product of the left defect
$a[bx]$ with the right defect $b[x]$, in this order
\cite[lines~4234--4245]{FrancoRubio2025MPUGauging}. Writing $e\leftrightarrow0$
and $g\leftrightarrow1$ for the labels, so that $bx=x+b$ in $\mathbb F_2$,
Eqs.~\ref{eq:fbc25_czx_defect_g1}--\ref{eq:fbc25_czx_defect_e} give
\begin{align}
    v(e,e,0)&=\ket{0000},
    &
    v(e,e,1)&=\ket{1111},
    \label{eq:fbc25_czx_v_ee}\\
    v(e,g,0)&=\ket{1100},
    &
    v(e,g,1)&=-i\ket{0011},
    \label{eq:fbc25_czx_v_eg}\\
    v(g,e,0)&=\ket{0000},
    &
    v(g,e,1)&=-i\ket{1111},
    \label{eq:fbc25_czx_v_ge}\\
    v(g,g,0)&=-i\ket{1100},
    &
    v(g,g,1)&=-i\ket{0011}.
    \label{eq:fbc25_czx_v_gg}
\end{align}
Each pair $v(a,b,0)$, $v(a,b,1)$ is orthogonal, as the modified-fusion
lemma requires. The four defect domains of
Eq.~\ref{eq:fbc25_czx_domain_definition} are
\begin{align}
    \mathcal K_{e,e}=\mathcal K_{g,e}
        &=\mathcal L_0
        :=\operatorname{span}\{\ket{0000},\ket{1111}\},
    \label{eq:fbc25_czx_even_domains}\\
    \mathcal K_{e,g}=\mathcal K_{g,g}
        &=\mathcal L_1
        :=\operatorname{span}\{\ket{1100},\ket{0011}\}.
    \label{eq:fbc25_czx_odd_domains}
\end{align}
Here $\mathcal L_p=\operatorname{span}\{v(e,p,x):x\}$ is the target support
of the product label $p$. The two supports are orthogonal,
\begin{align}
    \mathcal L_0\perp\mathcal L_1,
    \qquad
    \mathcal L_1\neq\mathcal L_0.
    \label{eq:fbc25_czx_targets_orthogonal}
\end{align}
The domain depends only on the right label $b$, because the single-site
defects $g[x]$ and $e[x]$ are proportional.

The prescribed maps $D_{a,b}\colon\mathcal K_{a,b}\to\mathcal L_{a+b}$,
$v(a,b,x)\mapsto v(e,a+b,x)$, of Eq.~\ref{eq:fbc25_czx_prescription} are
\begin{align}
    D_{e,e}&=\Id_{\mathcal L_0},
    &
    D_{e,g}&=\Id_{\mathcal L_1},
    \label{eq:fbc25_czx_maps_identity}\\
    D_{g,e}\ket{0000}&=\ket{1100},
    &
    D_{g,e}\ket{1111}&=\ket{0011},
    \label{eq:fbc25_czx_map_ge}\\
    D_{g,g}\ket{1100}&=i\ket{0000},
    &
    D_{g,g}\ket{0011}&=i\ket{1111}.
    \label{eq:fbc25_czx_map_gg}
\end{align}
Each is a unitary isomorphism between its domain and its target. The maps
$D_{g,e}$ and $D_{g,g}$ exchange the two supports, while $D_{e,e}$ and
$D_{e,g}$ preserve them.

\section{The displayed fusion operators restrict to the prescribed maps}

The movement operator $w_R$ of
Ref.~\cite[lines~4860--4888]{FrancoRubio2025MPUGauging} acts on the four
qubit lines with $Z$ on qubits $2,3$ at the bottom, a controlled-$Z$ between
qubits $1,2$, a controlled-$Z$ between qubits $0,1$, and $X$ on qubits $0,1$
at the top. The fusion operator $\lambda^R_{g,g}$ of
Ref.~\cite[lines~4989--5013]{FrancoRubio2025MPUGauging} carries the
prefactor $-i$, $X$ on qubit $2$ at the bottom, the same two controlled-$Z$
gates, and $X$ on qubits $0,1,2$ together with $Z$ on qubit $3$ at the top.
Read from bottom to top,
\begin{align}
    w
        &=(X_0X_1)\,\mathrm{CZ}_{01}\mathrm{CZ}_{12}\,(Z_2Z_3),
    \label{eq:fbc25_czx_w}\\
    \lambda
        &=-i\,(X_0X_1X_2Z_3)\,\mathrm{CZ}_{01}\mathrm{CZ}_{12}\,X_2,
    \label{eq:fbc25_czx_lambda}\\
    \tilde\lambda
        &=\lambda Z_0,
    \label{eq:fbc25_czx_tilde_lambda}
\end{align}
where Eq.~\ref{eq:fbc25_czx_tilde_lambda} is the modification
$\tilde\lambda^R_{g,g}=\lambda^R_{g,g}Z_1$ of
Ref.~\cite[lines~5179--5183]{FrancoRubio2025MPUGauging}. The source names
this factor ``the Pauli $Z$ operator on the first qubit'' and numbers the
qubits of a site from $1$, whereas the present note numbers the four qubits
of the two sites $0,1,2,3$. The source's $Z_1$ and the $Z_0$ of
Eq.~\ref{eq:fbc25_czx_tilde_lambda} are therefore one operator, the Pauli
$Z$ on the first qubit of the left site, written in the two numbering
conventions.\footnote{These are the gate forms
transcribed in \doclink{TNLean/MPS/MPDO/CZXGaussCircuitTuple}, whose
qubit numbering and phase tables the present note verifies against the
diagrams.} With $\overline x=(x_0+1,x_1+1,x_2,x_3)$, the phase tables are
\begin{align}
    w\ket{x}
        &=(-1)^{x_0x_1+x_1x_2+x_2+x_3}\ket{\overline x},
    \label{eq:fbc25_czx_w_phase}\\
    \lambda\ket{x}
        &=-i(-1)^{x_0x_1+x_1x_2+x_1+x_3}\ket{\overline x},
    \label{eq:fbc25_czx_lambda_phase}\\
    \tilde\lambda\ket{x}
        &=-i(-1)^{x_0x_1+x_1x_2+x_0+x_1+x_3}\ket{\overline x}.
    \label{eq:fbc25_czx_tilde_lambda_phase}
\end{align}

These transcriptions reproduce every fusion identity the source displays.
The two movement identities of
Ref.~\cite[lines~4889--4988]{FrancoRubio2025MPUGauging} read, in the
vectors of Eqs.~\ref{eq:fbc25_czx_v_eg} and~\ref{eq:fbc25_czx_v_ge},
\begin{align}
    w\,v(g,e,0)
        &=w\ket{0000}=\ket{1100}=v(e,g,0),
    \label{eq:fbc25_czx_w_identity_0}\\
    w\,v(g,e,1)
        &=-i\,w\ket{1111}=-i\ket{0011}=v(e,g,1),
    \label{eq:fbc25_czx_w_identity_1}
\end{align}
which is the statement $\lambda^R_{g,e}=w_R$ with trivial $L$-symbols
\papereqref{fbc25}{eq:simplecases} of
Ref.~\cite[lines~3331--3340]{FrancoRubio2025MPUGauging}. The two identities
displayed for $\lambda^R_{g,g}$ in
Ref.~\cite[lines~5014--5066]{FrancoRubio2025MPUGauging} read
\begin{align}
    \lambda\,v(g,g,0)
        &=-i\,\lambda\ket{1100}=-i\,(-i)\ket{0000}=-v(e,e,0),
    \label{eq:fbc25_czx_lambda_identity_0}\\
    \lambda\,v(g,g,1)
        &=-i\,\lambda\ket{0011}=-i\,(i)\ket{1111}=v(e,e,1),
    \label{eq:fbc25_czx_lambda_identity_1}
\end{align}
which are the printed $L$-symbols $L^0_{g,g}=-1$ and $L^1_{g,g}=1$
\cite[line~5067]{FrancoRubio2025MPUGauging}. The alternative top-to-bottom
reading of the circuits would give $w\,v(g,e,0)=-v(e,g,0)$ and
$\lambda\,v(g,g,1)=-v(e,e,1)$, contradicting the displayed identities; the
displayed identities therefore pin the reading. The modified operator then
satisfies
\begin{align}
    \tilde\lambda\,v(g,g,0)
        &=\lambda Z_0\bigl(-i\ket{1100}\bigr)
        =i\,\lambda\ket{1100}
        =\ket{0000}=v(e,e,0),
    \label{eq:fbc25_czx_tilde_identity_0}\\
    \tilde\lambda\,v(g,g,1)
        &=\lambda Z_0\bigl(-i\ket{0011}\bigr)
        =-i\,\lambda\ket{0011}
        =\ket{1111}=v(e,e,1),
    \label{eq:fbc25_czx_tilde_identity_1}
\end{align}
which is Eq.~\ref{eq:fbc25_czx_prescription} exactly. The $Z$ on qubit $1$
would do the same, because both vectors of $\mathcal L_1$ have equal first
two bits; a $Z$ on qubit $2$ or $3$, or the product $Z_0\lambda$, would
satisfy the prescription only up to a global sign $-1$. The source's
choice, $Z$ on the first qubit applied before $\lambda^R_{g,g}$, is the one
that satisfies Eq.~\ref{eq:fbc25_czx_prescription} without a sign.

Together with the trivial cases $\lambda^R_{e,e}=\lambda^R_{e,g}=\Id$,
Eqs.~\ref{eq:fbc25_czx_w_identity_0}--\ref{eq:fbc25_czx_tilde_identity_1}
show that the four displayed operators $(\Id,\Id,w,\tilde\lambda)$, indexed
by $(a,b)\in\{e,g\}^2$, restrict on the domains of
Eqs.~\ref{eq:fbc25_czx_even_domains} and~\ref{eq:fbc25_czx_odd_domains} to
the prescribed maps of
Eqs.~\ref{eq:fbc25_czx_maps_identity}--\ref{eq:fbc25_czx_map_gg}.

\section{Verdicts}

\paragraph{Convention matching is a theorem.}
The retained four-qubit circuits are the source's fusion operators in the
conventions of the preceding sections: qubits $0,1$ on the left blocked
site and $2,3$ on the right, circuits read from bottom to top, the
controlled-$Z$ decomposition of Eq.~\ref{eq:fbc25_czx_cz_decomposition},
the displayed prefactor $-i$, and the displayed fusion gauge
$\tilde\lambda^R_{g,g}=\lambda^R_{g,g}Z_1$. Every one of these conventions
is either stated by the source or forced by a displayed identity, and the
restrictions of the circuits to the derived domains are the prescribed maps.
No axiom is needed.

\paragraph{The single-image identification is false.}
The odd target support is
\begin{align}
    \mathcal L_1
        &=\operatorname{span}\{\ket{1100},\ket{0011}\}
        =\mathcal K_{g,g},
    \label{eq:fbc25_czx_L1}
\end{align}
which is orthogonal to $\mathcal L_0=\operatorname{span}\{\ket{0000},\ket{1111}\}$
by Eq.~\ref{eq:fbc25_czx_targets_orthogonal}. The source draws enough to
fix $\mathcal L_1$, so this is not a source gap. The identification
$\mathcal L_1=\mathcal L_0$ assumed in the planning material yields the
domains $\mathcal K_{e,g}=\mathcal L_0$ and $\mathcal K_{g,e}=\mathcal L_1$,
which are the wrong way round: by
Eqs.~\ref{eq:fbc25_czx_even_domains} and~\ref{eq:fbc25_czx_odd_domains},
$\mathcal K_{e,g}=\mathcal L_1$ and $\mathcal K_{g,e}=\mathcal L_0$. The
support-level classification by a Grassmannian of rank-two odd targets
collapses: the source fixes the odd target, and the four-sector
support-and-restriction system is unique. The support projector of the
sector $(a,b)$ is the orthogonal projector onto $\mathcal L_b$. Transports
between sectors with different product labels are the maps $D_{g,e}$ and
$D_{g,g}$ themselves, so no additional identification of the two supports is
required or available.

\paragraph{Membership in the full completion class.}
Because the displayed tuple $(\Id,\Id,w,\tilde\lambda)$ agrees with the
prescribed isometries on all four domains, it is a member of the full unitary
completion class of the CZX defect maps, the class of unitary tuples that
extend the four prescribed maps. The exact dimension $2^N$ of its
gauge-invariant subspace on $N\geq3$ blocked sites is therefore a lower
bound on the maximum of that dimension over the full class. The value $2^N$
is not stated in the source, whose outlook
\cite[line~5198]{FrancoRubio2025MPUGauging} leaves the growth of this
subspace to future work; it is a theorem of the formal development about the
displayed tuple alone, proved from the explicit gates.\footnote{The theorem is
\texttt{finrank\_commonFixedSubmodule\_placedGaussProjector\_circuitTuple}
in \doclink{TNLean/MPS/MPDO/CZXGaussInvariantSubspace}.} This lower bound
was previously conditional on supplying the three undisplayed domains and
maps; the derivation above supplies them.

\paragraph{An index convention that does not matter.}
The gauged tensor displayed in
Ref.~\cite[lines~5131--5178]{FrancoRubio2025MPUGauging} carries, for gauge
label $1$, the matrix with entries $1$ and $-i$ off the diagonal. With rows
indexed by the block to the right of the defect and columns by the block to
its left, its entries are the scalars of
Eqs.~\ref{eq:fbc25_czx_defect_g0} and~\ref{eq:fbc25_czx_defect_g1}; with
the opposite reading the two scalars are exchanged. Either reading gives the
same four domains and the same four prescribed maps, since the scalars enter
both sides of Eq.~\ref{eq:fbc25_czx_prescription} consistently.

\section{Formal boundary and elimination plan}

The single-site defects of
Eqs.~\ref{eq:fbc25_czx_defect_g1}--\ref{eq:fbc25_czx_defect_e}, the eight
two-site vectors of
Eqs.~\ref{eq:fbc25_czx_v_ee}--\ref{eq:fbc25_czx_v_gg}, the domains of
Eqs.~\ref{eq:fbc25_czx_even_domains} and~\ref{eq:fbc25_czx_odd_domains}, and
the prescribed maps of
Eqs.~\ref{eq:fbc25_czx_maps_identity}--\ref{eq:fbc25_czx_map_gg} are carried
by formal declarations, packaged as one family of prescribed defect maps
satisfying Eq.~\ref{eq:fbc25_czx_prescription} on their domains, together with
the orthogonality of Eq.~\ref{eq:fbc25_czx_targets_orthogonal} and the
two-dimensionality of each domain. The formal development also contains the
four circuits, their phase
tables, and the exact dimension of the gauge-invariant subspace of the
displayed tuple. The finite computation of
Eqs.~\ref{eq:fbc25_czx_w_identity_0}--\ref{eq:fbc25_czx_tilde_identity_1} is
also formal, so the displayed tuple is a member of the completion class of
those prescribed maps and the lower bound on the gauge-invariant dimension
over that class carries no conditional clause. What the formal development
does not contain is any statement about the minimum over the class, the exact
value of the supremum, or which completion attains it; these remain open as in
the source. The CZX defect and fusion data are
tracked by \ghissue{7355}, the prescribed defect maps by \ghissue{7739}, the
completion membership of the displayed tuple by \ghissue{7740}, the
completion-class instance by \ghissue{7569}, and the
gauge-invariant-subspace problem by \ghissue{7568}.

\bibliographystyle{alpha}
\bibliography{references}

\end{document}
